\section{DreamFusion：SDS 的代表性应用}

\subsection{DreamFusion 概述}

DreamFusion（Poole et al., 2022）是第一个将 SDS 成功应用于文本到 3D 生成的工作。它将预训练的 Imagen 文本到图像扩散模型的知识蒸馏到一个 NeRF 中，实现了仅从文本描述生成 3D 模型。

\begin{importantbox}{DreamFusion 的流程}
DreamFusion 的工作流程如下：
\begin{enumerate}
    \item 初始化一个随机的 NeRF 表示 $\theta$
    \item 在每次迭代中：
    \begin{itemize}
        \item 随机选择一个相机视角 $\pi$
        \item 渲染该视角下的 2D 图像 $x_0 = g(\theta, \pi)$
        \item 计算 SDS 梯度并更新 $\theta$
    \end{itemize}
    \item 经过数千次迭代后，$\theta$ 收敛为一个高质量的 3D 模型
\end{enumerate}
\end{importantbox}

\subsection{生成结果示例}

DreamFusion 能够根据各种创意性文本描述生成 3D 模型，包括：
\begin{itemize}
    \item ``a fox wearing a sweater''（穿毛衣的狐狸）
    \item ``a ghost eating a hamburger''（吃汉堡的幽灵）
    \item ``a DSLR photo of a peacock on a surfboard''（冲浪板上的孔雀）
\end{itemize}

这些都是在互联网上几乎找不到对应 3D 模型的创意场景，但由于图像扩散模型在数十亿张图像上训练过，它``见过''足够多的相关视觉概念，因此能够通过 SDS 引导生成合理的 3D 结果。

\subsection{SDS 的广泛应用}

SDS 不仅限于 3D 生成，它的框架可以推广到任何满足以下条件的任务：
\begin{enumerate}
    \item 目标数据可以通过某种\textbf{可微分映射}转换为图像
    \item 有对应的\textbf{预训练扩散模型}可用
\end{enumerate}

\begin{table}[H]
\centering
\begin{tabular}{lll}
\toprule
\textbf{应用场景} & \textbf{数据类型} & \textbf{渲染函数 $g$} \\
\midrule
文本到 3D & 3D NeRF/Mesh & 可微分 3D 渲染器 \\
矢量图生成 & SVG 笔画参数 & 可微分光栅化 \\
4D 生成 & 动态 3D 场景 & 视频渲染器 \\
3D 编辑 & 已有 3D 模型 & 可微分 3D 渲染器 \\
纹理生成 & 表面纹理贴图 & UV 映射渲染 \\
\bottomrule
\end{tabular}
\caption{SDS 框架在不同应用场景中的实例化}
\end{table}

\begin{knowledgebox}{从图像扩散到视频扩散}
SDS 的思想同样可以与预训练的\textbf{视频扩散模型}结合，用于生成 4D 内容（动态 3D 场景）。通过将视频帧作为渲染目标，可以利用视频扩散模型中学到的时间一致性先验来生成动态的 3D 场景。
\end{knowledgebox}

\subsection{本章小结}

DreamFusion 是 SDS 在文本到 3D 生成中的里程碑式应用，展示了将大规模 2D 先验迁移到 3D 领域的巨大潜力。SDS 的框架具有很强的通用性，可以推广到矢量图、4D 场景、纹理、音频等多种模态。

\newpage
